\documentclass{article}
\usepackage{graphicx} % Required for inserting images
\usepackage{amsmath}
\usepackage{algorithmic}
\usepackage{algorithm}
\usepackage[polish]{babel}
\usepackage[utf8]{inputenc}
\usepackage[T1]{fontenc}
\usepackage{float}
\usepackage{listings}
\usepackage[table]{xcolor}
\usepackage{diagbox}
\usepackage{geometry}

\title{Algorytmy Równoległe - Laboratorium 2}
\author{Adam Staniszewski}
\date{}

\begin{document}

\maketitle

\section*{Sito Eratostenesa}
Przygotojmy sito i sprawdźmy jakie wyniki osiągnie w wersji sekwencyjnej.

\begin{lstlisting}
def sieve_of_erathosthenes(limit: int) -> list[int]:
    sieve: list[bool] = np.ones(limit + 1, dtype=bool)
    sieve[0] = False
    sieve[1] = False

    for num in range(2, int(math.sqrt(limit)) + 1):
        if sieve[num]:
            sieve[num*num:limit+1:num] = False

    return np.nonzero(sieve)[0]
\end{lstlisting}
Znalezienie liczb pierwszych do 999999999 zajęło 9.84 sekund. Znaleźliśmy 50847534 liczb.

\section*{Zrównoleglenie sita}
Implementacja sita z poprzedniego zadania pozostanie bez zmiany, ale tym razem będzie ono wywoływane nie wprost, ale przez funkcję \textbf{parallel\_sieve}.

\begin{lstlisting}
def parallel_sieve(n):
    comm = MPI.COMM_WORLD
    rank = comm.Get_rank()
    size = comm.Get_size()

    sqrt_n: int = int(math.sqrt(n))

    # Punkt 1
    primes_in_B: list[int] = sieve_of_erathosthenes(sqrt_n)

    # Punkt 2
    start_C: int = sqrt_n + 1 + rank * (n - sqrt_n) // size
    end_C: int = sqrt_n + 1 + (rank + 1) * (n - sqrt_n) // size - 1

    # Punkt 3
    sieve: list[bool] = np.ones(end_C - start_C + 1, dtype=bool)

    for prime in primes_in_B:
        first_multiple = max(prime*prime, (start_C + prime - 1) // prime*prime)
        sieve[first_multiple - start_C:end_C - start_C + 1:prime] = False

    local_primes = np.nonzero(sieve)[0] + start_C

    # Punkt 4
    all_primes = comm.gather(local_primes, root = 0)

    if rank == 0:
        all_primes = np.concatenate([primes_in_B] + all_primes)
        return all_primes
    else:
        return None
 
\end{lstlisting}
Połączmy się z Aresem i uruchommy nasz algorytm równoległy na 1 węźle obliczeniowym z 4 rdzeniami (analogicznie do konfiguracji z pliku Quickstart). 

Algorytm równoległy znalazł liczby pierwsze do 999999999 w 4.97 sekundy. Podobnie do implementacji sekwencyjnej, znalazł ich 50847534. Widzimy więc, że przedział został poprawnie podzielony między procesy i nie zgubiliśmy żadnych wartości.

\section*{Eksperymenty}
Sprawdźmy jeszcze jak będzie zmieniał się czas działania algorytmu przy wzroście liczby rdzeni. Dla każdego wariantu uruchomimy algorytm 5 razy i wyciągniemy średni czas.

\begin{table}[ht]
\centering
\begin{tabular}{|l|r|r|r|}
\hline
\rowcolor{orange!40} 
Liczba rdzeni & Średni czas wykonywania algorytm  \\\hline
1 (implementacja sekwencyjna) & 9.63 sekund \\\hline
4 rdzenie &  4.87 sekund \\\hline
8 rdzeni &  3.89 sekundy \\\hline
16 rdzeni &  2.82 sekundy \\\hline

\end{tabular}
\caption{\label{tab:widgets}Średni czas obliczeń dla różnej ilości rdzeni.}
\end{table}

\section*{Wnioski}
Możemy zaobserwować, że zrównoleglenie i dostosowanie ilości rdzeni w weźle obliczeniowym może znacznie przyszpieszyć szukanie liczb pierwszych przy pomocy Sita Eratostenesa.

\end{document}
